\documentclass[11pt,letterpaper,english]{article}
\usepackage[T1]{fontenc} % Standard package for selecting font encodingsamely titles, headers and contents
\usepackage{txfonts} % makes spacing between characters space correctly
\usepackage{xcolor} % Driver-independent color extensions for LaTeX and pdfLaTeX.
%\usepackage[pagestyles,raggedright]{titlesec} % related with sections—n
%\usepackage{blindtext} % To create text
\usepackage{fancyhdr} % header footer placement

\usepackage[top=1in, bottom=1in, left=1in, right=1in] {geometry} % Margins
\usepackage{graphicx}  % Essential for adding images to you document.

\usepackage{sectsty}
\sectionfont{\normalsize}
\subsectionfont{\normalsize}
\subsubsectionfont{\normalsize \it}

\usepackage{caption}
\captionsetup{labelsep=period}


\pagestyle{fancy} % allows you to use the header and footer commands
\usepackage{ragged2e}

\raggedright
\begin{document}

\setlength{\parindent}{0in} % Amount of indentation at the first line of a paragraph.

\pagestyle{fancy} \lhead{Advancing Fusion and Fission Energy through Exascale: Year 3} \rhead{E. Merzari et al.} \renewcommand{%
%\pagestyle{fancy} \lhead{Harnessing Exascale: First of a Kind High Fidelity Full Core Simulations } \rhead{E. Merzari et al.} \renewcommand{%
%\pagestyle{fancy} \lhead{Large Eddy Simulation of Rod Bundle Cores} \rhead{E. Merzari et al.} \renewcommand{%
\headrulewidth}{0.0pt}

\begin{center}
\bf {PROJECT EXECUTIVE SUMMARY} \\
\end{center}




\bigskip

\textbf{Title}: Advancing Fusion and Fission Energy through Exascale: Year 3 \smallskip
%\textbf{Title}: High-Fidelity Full-Core Simulations \smallskip

\textbf{PI and Co-PI(s)}: Elia Merzari, Paul Fischer, Misun Min, April Novak, Jun Fang,  John Tramm, Patrick Shriwise, Paul Romano \smallskip

\textbf{Applying Institution/Organization}: Pennsylvania State University \smallskip

\textbf{Resource Name(s) and Number of Node Hours Requested}: \\
Frontier: 250,000 node-hours  \\
Aurora: 250,000 node-hours  \\
\smallskip

\textbf{Amount of Storage Requested}: 770 TB (Frontier), 900 TB (Aurora)  \smallskip

\textbf{Executive Summary}: \\
\justify

Advanced nuclear energy holds transformative potential as a reliable, carbon-free power source capable of meeting urgent climate and energy security goals. A global surge in investment in both fission and fusion technologies signals a pivotal transition from academic research to commercial viability. However, the design, certification, and licensing of novel reactor concepts remain formidable challenges, especially for fusion systems where integrated testing is costly and complex. This reality necessitates high-fidelity, multiphysics simulations to accelerate development and reduce deployment risks.

This research aims to provide the simulation capabilities essential for advancing fusion energy by enabling unprecedented insight into large-scale, multiphysics phenomena. Specifically, we will perform first-of-a-kind,  simulations of fusion reactor components; approaching integrated multiphysics modeling of breeder blankets and novel reticulated foam tritium extraction systems. These simulations will be executed on Frontier and Aurora, leveraging its exascale resources to push the boundaries of current modeling efforts.

Fusion reactor components present extreme challenges in thermal-fluid dynamics, magnetohydrodynamics (MHD), heat transfer, and mass transport, which cannot be captured by the reduced-scale or single-physics models that have dominated the nuclear landscape. Knowledge gaps in these phenomena currently limit predictive capabilities for key performance metrics such as temperature distribution, tritium breeding ratio, heat extraction efficiency, and component longevity. Addressing these gaps is critical to the ultimate success of commercial fusion systems.

This research represents a unique opportunity to deliver both scientific and strategic value, providing high-resolution datasets for deep physical understanding while enabling the development of closure models for rapid integration into industry-scale design tools. These insights will be foundational for guiding design optimization, safety validation, and licensing strategies in emerging fusion platforms.

We will employ NekRS, a GPU-native, highly scalable spectral element code derived from Nek5000 and recognized with Gordon Bell and R\&D 100 awards. NekRS has demonstrated excellent performance on over 9,000 nodes of Frontier and is ideally suited for the thermal-fluid challenges of fusion systems. For neutron and gamma transport, we will use OpenMC, a Monte Carlo code with proven performance on both Frontier and Aurora.

This project is timely and well-aligned with the mission of the Leadership Computing Facilities (LCFs) to enable scientific breakthroughs at scale. By executing these first-of-their-kind simulations, we will establish fusion energy research as a front-runner in exascale computing, setting the stage for accelerated commercialization of next-generation nuclear technologies.

\end{document}
